%% APPENDIX B -- the spatial screen: calibration, validation, and the one
%% worked example of the vetting chain.  Owner: appendix-triage.
%% Carries \label{app:radius} (cited by 04_method.tex) and
%% \label{sec:heldout} (cited from the other appendices).
\section{The spatial screen: calibration and validation}
\label{app:radius}

The spatial screen asks whether an excess sits at the star or anywhere in the
field, and it removes most of the unattributed crossings. Everything the
dispositions rest on therefore depends on how well that screen is
calibrated, which is what this appendix measures. The answer is adverse to
the screen and is stated as such: \rankstatus{}

The star lies at the phase centre, inside the control annulus, and that
asymmetry has a known instrumental mechanism. Images are corrected for the
primary-beam response, so the true noise rises with distance from the
pointing centre while the search divides every position by one scale
estimated over the whole window; the control probes therefore sit at
systematically different effective noise from the star. A radius-corrected
normalisation removes that term: each control probe is standardised within
its own radius bin, the probe radii being known from the fixed draw, with no
free parameter.

\emph{Applied out of sample, the correction does not repair the statistic.}
It was frozen --- probe radii, bin edges and standardisation alike --- and
applied unmodified to \RadValStarN{} windows that played no part in its
construction. The uncorrected stellar rank has median \RadValStarRawMed{}
against the exchangeable value of 0.5; after correction the median is
\RadValStarCorMed, moving \RadValStarToward{} it, and the result holds
separately in both out-of-sample partitions (hold-out
\RadValStarHoRawMed$\to$\RadValStarHoCorMed, extension
\RadValStarExtRawMed$\to$\RadValStarExtCorMed). Rescoring the whole survey
with the gradient removed takes the rank-passing list from
\LocStageGlobal{} windows to \LocStageLocal, losing \LocStageLost{} and
gaining \LocStageGained, all marginal and all dispositioned on other
grounds, and the median rank over the \LocAllWin{} windows that retain a
control vector moves from \LocAllMedGlobal{} to \LocAllMedLocal, slightly
further from 0.5.

The interpretation is narrow and, we think, unavoidable. The radial gradient
is real and is one mechanism behind the displacement, but removing it does
not leave an exchangeable statistic, so the residual is not radial. Neither
normalisation yields a calibrated tail probability. \textbf{We therefore
retain the frozen screen as primary --- it has at least the virtue of having
been fixed in advance --- use it only to decide which windows are examined,
and rest every disposition on physical evidence or on recurrence rather than
on either rank.} The radial correction, the region-maximum audit and the
pseudo-star tests each leave every disposition unchanged.

\subsection{The reserved hold-out, and what it calibrates}
\label{sec:heldout}

The statistic and the mask were settled on the searched sample and cannot
also calibrate themselves, so a hold-out was reserved by a rule fixed and
recorded before the archival search finished and before any reserved block
was searched: a block is held out if and only if
$\mathrm{sha256}(\textsc{uid})_{[:8]} \bmod 5 = 0$, subject to two guards
that use identifiers alone --- a block on the pre-registered block list
deposited with this paper stays in the survey, and no system loses every
block --- so that the assignment can be recomputed from the UIDs and that
list alone. The three choices that can manufacture a
candidate were fixed before the hold-out rule itself: the detection statistic
of Eq.~\ref{eq:tstar}, the promote and retire criteria, and then the
reservation, with the first reserved block searched \HoLeadBlockDays{} days
later. The diagnostics built afterwards --- the radius-corrected
normalisation, the visibility-domain test, the drift-resolved injection
analysis and the $P_{90}$ budget --- cannot promote a window.

Of \NEbProcessedAll{} blocks processed, \NEbHeldOut{} are reserved
(\HoPctHeld{} per cent), giving \HoWin{} windows toward \HoStars{} stars with
no star and no system lost from the headline sample. Three quantities are
calibrated on those blocks and on nothing else.

First, the rank displacement. The stellar add-one rank out of sample has
median \HoRankMed{} where 0.5 is expected, with a block-clustered 95 per cent
interval of \HoRankMedLo--\HoRankMedHi{} over \HoWin{} windows, and
\HOFracLoPct{} per cent of ranks below 0.1 against \HOFracHiPct{} per cent
above 0.9, so the distribution is displaced rather than growing a tail.
Clustering removes neither the shift nor its size: \HOBlockBelow{} of
\HOBlockN{} blocks have a mean rank below 0.5, and a star-clustered bootstrap
places the median at \HOBootMed{} (95 per cent \HOBootLo--\HOBootHi). The
searched sample leans the same way without resolving it (median \SurvMedP,
star-clustered 95 per cent \SurvBootLo--\SurvBootHi), and the two
distributions are indistinguishable, so the in-sample checks were
underpowered rather than contradicted. A displacement that survives a
hold-out chosen before the data were seen is a property of the statistic and
not of this sample, which is the quantitative reason the rank is a screen and
never a significance.

Second, the false-alarm tail factor --- the measured rate at which a
no-signal position ranks first, divided by the rate exchangeability would
give. It is \HoTailFactor{} on \HoPseudoTrials{} trials, with a bootstrap
clustered over the \HoTailBlocks{} reserved blocks giving
\HoTailLo--\HoTailHi. That factor and its width set the accuracy of every
chance expectation in this paper. The reserved blocks contain \HoStageOne{}
rank-passing crossings against \HoStageOneExp{} expected.

Third, the radial behaviour of the control ensemble. The standardised
control level runs \HoRadBins{} from the inner to the outer edge of the
annulus, a mean per-window slope of $\HoRadSlope\pm\HoRadSlopeSe$: the
displacement reproduces out of sample, but the steep inner-edge excess seen
in sample does not --- which is why the radial term is one demonstrated
mechanism for the displacement rather than the whole of it.

A second out-of-sample set is used in this paper and must not be conflated
with the hold-out. The \emph{external calibration sample} is \CampBlocks{}
execution blocks toward \CampStars{} stars, reserved by a rule fixed before
the archival search finished, and it is not a subset of the in-scope blocks
of \S\ref{sec:sample}: it includes blocks toward stars outside the 40\,pc
work list, so its block count cannot be placed inside that ledger and enters
no count of the science sample. It is used only to calibrate the rank
statistic, and it carries the repeat observations the hold-out does not. Over
its \CampWindows{} windows the frozen pipeline returns \CampStageOne{}
rank-passing crossings, of which \CampStageOneCO{} are $\beta$~Pictoris CO
--- the positive control recovering on data it had never seen --- and
\CampUnattrib{} are unattributed, against \CampExpUn{} expected at the
measured tail rate.

\subsection{One crossing through the whole chain}
%% (subsection label retired: nothing cites it)

CP$-$72~2713 is a worked example of the vetting chain rather than a result,
and it is the clearest one the survey contains because every stage returns a
verdict. At the chance expectation of Appendix~\ref{app:falsealarm},
\ClustStageMean{} such events are expected over the Class~A windows and
\NStageOneUnattrib{} are observed.

In Band~7 the crossing reaches $T_\star=\VnCpT{}$ and outranks every one of
the \NCtrl{} control positions, so it passes the spatial screen; it is
channel \VnCpChan{} of its window; and no transition of the frozen mask lies
within $\pm\MaskHalfKms$\,km\,s$^{-1}$ of it, so it is unattributed.
CO($3{\to}2$) falls in this same window and is absent, below a $5\sigma$
limit fainter than the crossing itself, and sulphur-bearing species are not
in the mask and were not in it before the search. The drift-following
visibility fit places the emission at the stellar position,
${\rm Re}/\sigma=\VnCpReAdopted$ with ${\rm Im}/\sigma=\VnCpImAdopted$
against a control-position maximum of \VnCpCtrlAdopted; by the argument of
\S\ref{sec:vistest} that is a consistency check and not a promotion.

\emph{What disposes of it is the repeat block.} That second block matches the
first in tuning, channel width, channel count, integration count and
on-source time, registers to \CpRecOffChan{} of a channel, and is the deeper
of the two, so a persistent emitter at the first block's flux would have
returned $T_\star=\CpRecExpT$. Pinned to the first block's own frequency and
drift it returns \CpRecT, and read as two measurements of one constant
amplitude the pair is inconsistent at $\CpRecExclPair\sigma$. The blocks are
\CpSepH\,h apart, so this is a within-night test: it excludes a persistent
emitter of that strength on that baseline and excludes nothing about an
intermittent one. \textbf{The disposition rests on that recurrence test and
not on the rank}, which is the chain this paper applies to every crossing it
cannot attribute.

%% ------------------------------------------------------------------
%% REQUEST (appendix-triage -> results, 05_results.tex):
%%   CP-72 2713 appears nowhere in the body.  It is one of the two
%%   unattributed crossings that PASS the spatial screen, and the ledger's
%%   own chain string for it is "unattributed; rank-flagged; tested and does
%%   not recur".  It is quoted here as a worked example only.  If the results
%%   section means to name the two rank-passing crossings -- and it should --
%%   this is one of them, and the recurrence exclusion at \CpRecExclPair
%%   sigma is the result, not the rank.
%% NOTE (pre-existing contradiction, now fixed here):
%%   this paragraph previously gave T_star = 5.81 against "a ring maximum of
%%   5.68" and then concluded the crossing "still fails the rank screen".
%%   5.81 < 5.68 is false and so was the conclusion: ledger_v403.csv has
%%   rank_screen = True for this row.  The ring maximum had no macro, so the
%%   statement is now made in the form the screen actually tests -- the star
%%   outranks all \NCtrl controls -- and the six typed literals in the
%%   paragraph (5.81, 5.68, the SO offset, the transitions per GHz, the
%%   91 per cent coincidence rate and the 3-against-2 channel width) are
%%   gone: \VnCpT replaces the first, and the sentences carrying the others
%%   are deleted because no macro exists for them.
%% REQUEST (appendix-triage -> numbers):
%%   If the SO line-density argument is wanted back, it needs three macros:
%%   transitions per GHz over this window, the coincidence probability at
%%   +/- \MaskHalfKms km/s, and the feature's recovered width in channels
%%   against the instrumental response.  Without them it cannot be typeset.
%% ------------------------------------------------------------------
